\begin{helper}[Hereditary WQO and strict well-foundedness]
Let $Q$ be a type, let $r$ be a quasi-order on $Q$, and let
$\mathsf{HerCtblPower}(Q)$ be the hereditary countable carrier from
\noderef{HerCtblPower}, with relation $\mathsf{HerCtblRel}(r)$ from
\noderef{HerCtblRel}.  Then
\[
\mathsf{WellQuasiOrdered}(\mathsf{HerCtblRel}(r))
\quad\Longleftrightarrow\quad
\mathsf{WellFounded}\!\left(\lambda a\,b.
  \mathsf{HerCtblRel}(r)(a,b)\mathbin{\wedge}
  \neg\mathsf{HerCtblRel}(r)(b,a)\right).
\]
The strict relation is oriented from the later, strictly smaller element to
the earlier element in a descending chain.
\end{helper}
\begin{proof}
For the forward implication, \noderef{HerCtblPowerIsPreorder} gives that
$\mathsf{HerCtblRel}(r)$ is a preorder on the hereditary carrier.  The
standard strict-part theorem for a well-quasi-order therefore gives
\[
\mathsf{WellFounded}\!\left(\lambda a\,b.
  \mathsf{HerCtblRel}(r)(a,b)\mathbin{\wedge}
  \neg\mathsf{HerCtblRel}(r)(b,a)\right),
\]
which is the required conclusion.

For the converse, argue by contradiction.  If the hereditary carrier were
not well-quasi-ordered, the sequence characterization of
$\mathsf{WellQuasiOrdered}$ would give a sequence
$a:\mathbb N\to\mathsf{HerCtblPower}(Q)$ such that
\[
 m<n\quad\Longrightarrow\quad
 \neg\mathsf{HerCtblRel}(r)(a(m),a(n)).
\]
For each $n$, form the hereditary tail
\[
 T(n)=\left\langle
   \mathsf{node}\bigl(\mathsf{ULift}(\mathbb N),
      k\mapsto (a(n+k.\mathord{.down}))._1\bigr),
   \mathsf{HereditarilyCountable}\text{-witness}
 \right\rangle.
\]
The index type $\mathsf{ULift}(\mathbb N)$ is nonempty and countable.  The
first component uses the tagged full-hierarchy node constructor from
\noderef{PowerQ}, and the witness in the second component is supplied by
\noderef{HerCtblNodeClosure}; hence this pair is an element of the subtype
$\mathsf{HerCtblPower}(Q)$.

We claim that for every $n$,
\[
 \mathsf{HerCtblRel}(r)(T(n+1),T(n))
 \quad\text{and}\quad
 \neg\mathsf{HerCtblRel}(r)(T(n),T(n+1)).
\]
By \noderef{HerCtblRel}, the first relation is the corresponding
$\mathsf{PowerRel}(r)$ comparison between the two node presentations.  By
the node--node reduction \noderef{PowerRelNodeNode}, it suffices for each
left index $k$ to choose the right index $k+1$.  The two selected children
are both the underlying presentation $a(n+1+k.\mathord{.down})$, so their
comparison is reflexive by \noderef{PowerRelRefl}; the arithmetic identity
$n+1+k.\mathord{.down}=n+(k+1).\mathord{.down}$ identifies the two displayed
children.

For the reverse relation, suppose that
$\mathsf{HerCtblRel}(r)(T(n),T(n+1))$.  Applying
\noderef{PowerRelNodeNode} to the zeroth left index supplies some index $k$
with
\[
 \mathsf{HerCtblRel}(r)(a(n),a(n+1+k.\mathord{.down})).
\]
But $n<n+1+k.\mathord{.down}$, so this contradicts the badness of $a$.
Thus the displayed strict relation holds at every $n$.

Consequently $T$ is an infinite descending chain for the relation
\[
 (x,y)\longmapsto
 \mathsf{HerCtblRel}(r)(x,y)\mathbin{\wedge}
 \neg\mathsf{HerCtblRel}(r)(y,x),
\]
because its successive pair is $(T(n+1),T(n))$.  This contradicts the
assumed well-foundedness by
$\mathsf{wellFounded\_iff\_isEmpty\_descending\_chain}$.  Hence the
hereditary carrier is well-quasi-ordered, proving the converse.
\end{proof}
